\section{Invocation and Configuration}
\label{sec:invocation}

\subsection{Command line}

\begin{req}{CLI-1}
The command is \code{fsb [flags] [path]}.  With no flags and no path it serves the
user's home folder.
\end{req}

\begin{req}{CLI-2}
The optional \code{path} names the first root.  More than one path argument is an
error.  A path that does not exist, or is not a folder, is an error.  Symbolic links
in a root are resolved, so a root is always a real location.
\end{req}

\begin{req}{CLI-3}
\code{--root PATH}, which may be repeated, adds further roots (for example
\code{/Volumes/Data}).  All roots are equal: each is listed as a place the user can
switch to when there are several.
\end{req}

\begin{req}{CLI-4}
The folder \code{/} is refused as a root unless \code{--allow-system-root} is given;
the message says so.  Rules (Section~\ref{sec:rules}) apply to a root of \code{/} as
to any other root.
\end{req}

\begin{req}{CLI-5}
By default the launch URL is opened in the user's default browser.  \code{--no-open}
only prints it.  \code{--browser NAME} opens it in the named macOS application
instead (for example \code{"Google Chrome"}).  A name that begins with \code{-} or
contains a NUL, carriage return or newline is refused before anything starts.  If
the browser cannot be opened, \texttt{fsb} says so, keeps running, and the printed
URL remains usable.  On systems other than macOS the URL is only printed.
\end{req}

\begin{req}{CLI-6}
\code{--debug} enables verbose logging on standard error and makes a request for a
denied path answer with a body that names the matching rule (status still 404).  No
other flag or setting reveals which rule matched.
\end{req}

\begin{req}{CLI-7}
\code{--init} writes the default rule files (Section~\ref{sec:config}) and exits.  It
creates the folder \code{\~{}/.config/fsb} with mode \code{0700} if needed, creates each
file with mode \code{0600}, and \emph{never overwrites} a file that exists: it reports
that the file was left unchanged.  It is the only operation in which \texttt{fsb}
creates a file.
\end{req}

\begin{req}{CLI-8}
\code{-h} or \code{--help} prints the usage text and exits with status 0.  An
unknown flag prints the error and the usage text and exits with status 1.
\end{req}

\subsection{Startup output}

\begin{req}{CLI-9}
On success \texttt{fsb} prints, on standard output, one line naming the roots it
serves and the loopback address and port, then one line giving the launch URL.  The
port is chosen by the operating system at each start and is never configurable.
\end{req}

\begin{req}{CLI-10}
The launch URL has the form \code{http://127.0.0.1:PORT/PREFIX/?token=TOKEN}, where
\code{PREFIX} and \code{TOKEN} are random values generated at each start
(Section~\ref{sec:security}).  It works once; opening it a second time is refused.
\end{req}

\begin{req}{CLI-11}
If one or more core deny rules are not in effect (\rid{RUL-16}),
\texttt{fsb} prints a warning on standard error at startup naming the paths that are
therefore browsable.
\end{req}

\subsection{Running and stopping}

\begin{req}{CLI-12}
\texttt{fsb} runs in the foreground until it receives \code{SIGINT} (Control-C) or
\code{SIGTERM}.  It then stops accepting requests, allows requests in progress up to
three seconds to finish, and exits with status 0.
\end{req}

\begin{req}{CLI-13}
The exit status is 0 after \code{--help}, \code{--init} and a normal shutdown, and 1
after any error.  An error is reported on standard error as one line beginning
\code{fsb:}, and no server is started.
\end{req}

\begin{req}{CLI-14}
\texttt{fsb} refuses to start, with a message that names the file and the line, if
the deny file or the ignore file cannot be parsed
(Section~\ref{sec:rulefiles}).  It never starts with weaker rules than the user wrote.
\end{req}

\subsection{Environment}

\begin{req}{ENV-1}
The home folder is taken from the operating system's notion of the user's home
(the \code{HOME} environment variable), then resolved to its real location, so that
a home folder reached through a symbolic link or an alias is treated as the real
folder.  It determines the default root, the location of the configuration folder,
and the meaning of \code{\~{}} in rules.
\end{req}

\subsection{Configuration files}
\label{sec:config}

\begin{req}{CFG-1}
Configuration is two optional text files in \code{\~{}/.config/fsb/}: \code{ignore}
(hide rules) and \code{deny} (block rules).  Both are global; there are no
per-folder rule files.
\end{req}

\begin{req}{CFG-2}
The files are read once, at startup.  Changes take effect at the next start.
\end{req}

\begin{req}{CFG-3}
If \code{ignore} does not exist, the built-in ignore rules apply (\code{.DS\_Store},
\code{.git/}, \code{node\_modules/}).  If \code{deny} does not exist, the core deny
rules of Section~\ref{sec:core} apply.  If \code{deny} exists it is authoritative:
whatever it does not contain is not enforced, and a core rule that is missing is
reported as in \rid{CLI-11} and \rid{RUL-14}.
\end{req}

\begin{req}{CFG-4}
User preferences (which columns are shown, the width of the preview pane, and the
like) are not configuration files: they are kept in the browser
(Section~\ref{sec:prefs}) and never on disk.
\end{req}
